\chapter{Experimental Design and Reproduction}\label{ch:method}

\section{Data sources and original computational design}
The principal inputs are the supplied eight-page research manuscript, its 42-page book-length draft, the seven-page preliminary report, the original \code{ulp.ipynb} notebook, four Plotly HTML exports, and the associated source-code archive. The preliminary report investigates $c=2^{-k}$, $k=1,\ldots,20$, and $c$ from two further families: odd integers and reciprocals of odd denominators. The manuscript and book draft extend these investigations to a systematic precision-dependent model.

The supplied validation code specifies three types of $c$, for each of binary64 and binary32: 1,001 odd integers from $1$ to $2,001$; 300 generated dyadic fractions; and 310 non-dyadic inputs consisting of ten named constants or fractions and 300 sampled real inputs. These categories total \textbf{1,611 increments per precision}. The stated random generator is \code{numpy.random.default_rng(42)}. The code samples 1,500--4,000 values of $x$ per relevant binade, computes $y=x^{10}$, stores $y$ once, and evaluates $(y+c)-y$. Results are grouped using the actual exponent of the computed $y$, with samples for which $y+c$ crosses a binade border handled separately.

\section{Floating-point conventions}
Arithmetic results are compared at the precision used for each experiment. The binary64 experiments use NumPy's \code{float64}, while binary32 experiments cast both $x$ and $c$ to \code{float32} before arithmetic to avoid unintentional promotion. Given a finite value $y$, the upward ULP in the numerical experiments is obtained from
\begin{equation}\label{eq:nextafter}
 \ulp(y)=\code{nextafter}(y,+\infty)-y.
\end{equation}
For every sample the recorded quantities are the increment $c$, the actual $y$, the result $f$, its relative error $(f-c)/c$, and (where applicable) its ULP-scaled error $(f-c)/\ulp(y)$.

\section{Focused replication and figure generation}
A smaller deterministic check was independently executed during report preparation: $x$ takes the 39 integers from $2$ through $40$, while $k$ takes the 20 integers from $1$ through $20$, with $c=2^{-k}$ and binary64 arithmetic. The 780 results are classified as preserved, absorbed, tied downward, or tied upward. The resulting counts match the supplied preliminary report and are given in Table~\ref{tab:powercount}.

All included visualizations are generated by \code{scripts/generate\_figures.py}, which performs the arithmetic again with NumPy and saves static PNG plots with English axes and captions. The power-of-two map, continuous curves, odd-integer curves, and odd-reciprocal curves are counterparts of the four supplied interactive HTML exports. Additional plots recreate the ULP staircase, precision comparison, odd-integer error map, and the model's normalized rounding pattern described in the research manuscript. The original interactive HTML is not directly embedded in the PDF, since its interface labels and tooltips are in Portuguese and interactive JavaScript is not portable in a conventional LaTeX PDF.

\section{Reproducibility and evidential status}
The figures and the 780-case enumeration were regenerated locally. The 1,611-per-precision validation outcome, large-scale binade-border frequencies, and the original source tables are reported from the supplied study and were \emph{not} rerun at their full sample sizes in this document. This distinction is important because the relevant conditions concern rounding behavior at extremely sparse boundary points. Source scripts are included in the accompanying materials for checking those experiments and for adapting them to new data.
